%% =========================================================================
\section{Model Limitations}
%% =========================================================================

Several simplifications in the Simscape Electrical model should be noted when
interpreting the results. The transmission lines between the generators and the
load bus are modelled as ideal (zero-impedance) connections, so line losses
and voltage drops along the network are absent. The synchronisation breaker
closes at a fixed time rather than waiting for the synchroscope null, which
introduces transient currents that would be smaller in a real synchronisation
procedure. Saturation effects in the machine's magnetic circuit are not
modelled; at high excitation levels, a real generator would exhibit reduced
flux linkage and altered transient response. Finally, the governor and turbine
are represented by first-order transfer functions, omitting boiler dynamics,
valve travel limits, and dead-band non-linearities present in physical
steam-turbine plant. These simplifications are standard for a power-system
stability study at this level, but they mean the quantitative results
(e.g.\ exact settling times and overshoot magnitudes) should be treated as
indicative rather than precise predictions of real-world behaviour.
